\section{Taxonomy}

The track is organized around five roles: stochastic search and ensembles,
exact optimization, multi-objective selection, legal posture, and audit
certification. The shared RPLAN/RCTX fixed point is the artifact boundary that
lets heterogeneous methods be compared without pretending they prove the same
kind of claim.

\begin{table}[htbp]
\centering
\small
\begin{tabular}{p{0.18\linewidth}p{0.34\linewidth}p{0.34\linewidth}}
\toprule
Paper band & Role & Primary claim type \\
\midrule
U.1--U.11 & Existing search and ensemble foundations & Method and empirical framing \\
U.12 & Algorithm-selection matrix & Methodology and operator guidance \\
U.13 & Exact-vs-heuristic certification & Certificate interpretation \\
U.14 & Multi-objective selection & Selection and export discipline \\
U.15 & Legal postures for search & Legal/policy interpretation limits \\
U.16--U.19 & Exact and heuristic implementations & Implementation and evaluation \\
U.20 & Plan audit certificates & Artifact and verifier contract \\
\bottomrule
\end{tabular}
\caption{U-series spine after the RPLAN/RCTX fixed point.}
\end{table}

The taxonomy is deliberately role based. A branch-and-cut run, a large
neighborhood search run, and a Pareto-front selection run may all emit usable
plans, but their evidence differs: one may offer a bound or infeasibility
witness, another may offer a validity-preserving improvement trace, and another
may offer a documented trade-off choice among non-dominated alternatives.
